\documentclass[11pt]{article}
\usepackage[margin=2cm]{geometry}
\usepackage[T1]{fontenc}
\usepackage{lmodern}
\usepackage{microtype}
\usepackage{amsmath,amssymb}
\usepackage{booktabs,tabularx}
\usepackage{enumitem}
\usepackage[hidelinks]{hyperref}
\setlist{nosep,leftmargin=*}
\setlength{\parindent}{0pt}
\setlength{\parskip}{4pt}

\title{\vspace{-1.5em}\textbf{Pruning the Air}\\[2pt]
\large Why Attention-Guided Visual Token Pruning Fails on Medical Images}
\author{Pushkar et al.}
\date{Research brief for CVPR 2027 \,$\cdot$\, 31 October 2026}

\begin{document}
\maketitle
\vspace{-1.5em}

% ------------------------------------------------------------------
\section*{0\quad Thesis}
\noindent\fbox{\parbox{\dimexpr\linewidth-2\fboxsep-2\fboxrule\relax}{%
On medical images, the attention signal that FastV-style pruners use to rank visual tokens
selects \emph{acquisition background} (air, empty field of view) rather than anatomy. At
aggressive budgets this makes the selection \textbf{worse than random}. Removing background
before scoring should recover most of the loss, training-free and at negligible cost.}}

{\small\emph{Status: hypothesis grounded in published results; our experiments start 2 Nov.}}

% ------------------------------------------------------------------
\section*{1\quad Problem}
\begin{itemize}
  \item Medical MLLMs (Lingshu-7B, HuatuoGPT-Vision-7B~\cite{huatuo}) pass hundreds to thousands
        of visual tokens to the LLM; prefill dominates inference cost on clinical hardware.
  \item The standard remedy is training-free visual token pruning. The dominant family ranks
        tokens by the LLM's own text-to-visual attention: FastV~\cite{fastv} (prunes at layer
        $K{=}2$), PDrop~\cite{pdrop}.
  \item FastV assumes \emph{high attention $\Rightarrow$ important token}, validated only on
        natural images at ${\sim}50\%$ retention. On medical images this assumption fails
        badly enough to fall below chance, and no work explains why.
\end{itemize}

% ------------------------------------------------------------------
\section*{2\quad Evidence (published results)}
\begin{center}\small
\begin{tabular}{@{}llrrrr@{}}
\toprule
Backbone & Benchmark (modality) & Full & Random & FastV & $\Delta$ \\
\midrule
Lingshu-7B  & OmniMedVQA (multi-modality)       & 80.62 & 38.69 & 32.26 & $\mathbf{-6.43}$  \\
            & SLAKE (CT/MRI/X-ray)              & 83.92 & 56.93 & 55.01 & $-1.92$           \\
            & RadImageNet-VQA, closed (CT/MRI)  & 89.40 & 55.59 & 37.50 & $\mathbf{-18.09}$ \\
            & RadImageNet-VQA, open (CT/MRI)    & 70.25 & 43.05 & 28.70 & $\mathbf{-14.35}$ \\
\midrule
Qwen3-VL-8B & OmniMedVQA (multi-modality)       & 68.16 & 49.92 & 51.49 & $+1.57$           \\
            & SLAKE (CT/MRI/X-ray)              & 80.08 & 49.25 & 66.92 & $+17.67$          \\
            & RadImageNet-VQA, closed (CT/MRI)  & 89.70 & 64.32 & 52.10 & $\mathbf{-12.22}$ \\
            & RadImageNet-VQA, open (CT/MRI)    & 73.65 & 47.05 & 37.70 & $\mathbf{-9.35}$  \\
\bottomrule
\end{tabular}

\vspace{2pt}
{\footnotesize From MedSTP~\cite{medstp}, 10\% of visual tokens kept. $\Delta=$ FastV $-$ Random.
Bold: FastV significantly below Random ($|\Delta|>2\,\mathrm{SE}$, from the reported bootstrap SEs).}
\end{center}

\begin{itemize}
  \item \textbf{Mild pruning too.} ViToS~\cite{vitos} (ICML'26), 80\% kept: FastV is the lowest
        of four training-free pruners on Lingshu-7B (avg.\ 60.3 vs.\ 62.9--63.3; full 63.7) and
        HuatuoGPT-Vision-7B (55.3 vs.\ 55.5--56.8; full 56.8). On OmniMedVQA it alone loses
        3.8 points; VisionZip, VScan and PDrop lose $\le 0.3$.
\end{itemize}

% ------------------------------------------------------------------
\section*{3\quad What prior work shows (natural images only)}
\begin{itemize}
  \item \textbf{Positional bias}~\cite{feather} (ICCV'25): early-layer attention keeps the bottom
        of the image (mean kept position 80.7\% down); on localization FastV scores 5.9 vs.\ 28.0
        for uniform sampling (unpruned 53.2). Attributed to 1D-RoPE decay.
  \item \textbf{Padding/sink inflation}~\cite{pore}: attention favors late and semantically
        empty (padding) tokens.
  \item \textbf{Random parity at depth}~\cite{randtok}: in deep layers pruning $\approx$ random,
        attributed to vanishing visual information.
  \item \textbf{Artifact tokens}~\cite{registers,ttreg}: ViTs create high-norm tokens on uniform
        background patches that absorb attention; removable at test time.
  \item \textbf{Benchmarks too easy}~\cite{vtcbench}: plain downsampling beats many compression
        methods.
\end{itemize}
\textbf{Not covered:} medical images; failure on ordinary VQA (not localization); M-RoPE
backbones (Qwen2.5/3-VL, the base of Lingshu); acquisition background as the trigger. Medical
methods propose new scorers or trained policies without diagnosing the failure
(MedSTP~\cite{medstp}, ViToS~\cite{vitos}).

% ------------------------------------------------------------------
\section*{4\quad What the published numbers already tell us}
\emph{Re-analysis of \S2; no experiments of ours yet.}
\begin{enumerate}[label=\textbf{D\arabic*}]
  \item \textbf{Worse than chance.} FastV is significantly below Random in 5/8 cells (by
        6.4--18.1 points); 2 ties, 1 above.
  \item \textbf{Modality-dependent.} Only RadImageNet-VQA (CT/MRI only) is below Random on
        \emph{both} backbones, with the largest gaps; multi-modality sets are inconsistent.
  \item \textbf{Not an extreme-budget artifact.} The FastV-specific loss is already visible at
        80\% kept.
  \item \textbf{Unexplained by prior accounts.} \cite{feather} covers 1D RoPE and localization;
        here the backbones use M-RoPE and the task is ordinary VQA. \cite{randtok} covers deep
        layers; FastV prunes at layer 2. PDrop (attention-based, multi-stage, later layers)
        loses 0.4 on Lingshu-7B vs.\ 3.4 for FastV, suggesting the weak point is \emph{early
        single-shot} attention.
\end{enumerate}

% ------------------------------------------------------------------
\section*{5\quad Hypotheses (each falsifiable)}
\begin{center}\small
\begin{tabularx}{\linewidth}{@{}l>{\raggedright}X>{\raggedright}X>{\raggedright\arraybackslash}X@{}}
\toprule
   & Mechanism & Prediction & Rejected if \\
\midrule
H1 & \textbf{Background capture.} Air in CT/MRI behaves like padding; layer-2 attention
     over-selects it.
   & Background share of FastV-kept tokens $>$ image background share ($=$ Random's). The
     Random$-$FastV gap grows with background fraction: CT/MRI $>$ X-ray $>$ fundus $>$
     derm./path.\ $\approx 0$.
   & FastV background retention $\le$ Random's. \\
\addlinespace
H2 & \textbf{Positional bias under M-RoPE.} Kept positions skew to non-anatomical regions.
   & Pattern differs between 1D-RoPE (HuatuoGPT-Vision) and M-RoPE (Lingshu); gap shrinks
     as $K{:}\ 2\!\to\!8\!\to\!16$.
   & Kept positions uniform ($\chi^2$ n.s.). \\
\addlinespace
H3 & \textbf{Artifact sinks.} High-norm visual tokens on uniform background attract the
     scoring attention.
   & Such tokens are selected far above their base rate.
   & Selected at base rate. \\
\addlinespace
H4 & \textbf{Small evidence} (residual).
   & After removing H1--H3, loss grows as evidence area shrinks (SLAKE organ masks).
   & Flat across evidence-size strata. \\
\bottomrule
\end{tabularx}
\end{center}

% ------------------------------------------------------------------
\section*{6\quad Proposed fix: Foreground Gate (training-free, zero parameters)}
Body mask $M=\mathrm{Dilate}\big(\mathrm{FillHoles}(\mathrm{LCC}(\neg B(x)))\big)$, with $B$ the near-black (or near-white) pixels connected to the image edge: the
largest connected component drops the scanner table; hole filling keeps the lungs, which are
as dark as outside air. Let $m_i\in[0,1]$ be the fraction of token $i$'s patch inside $M$, and
$s_i$ any base score (e.g.\ FastV attention at layer $K$). With budget $k=\lfloor\rho N\rfloor$:
\[
\tilde s_i=\begin{cases} s_i, & m_i\ge\tau\\ -\infty, & \text{otherwise}\end{cases}
\qquad
\mathcal{K}=\operatorname{TopK}_k(\tilde s)
\]
(if fewer than $k$ tokens pass the gate, the remainder is filled uniformly).
\begin{itemize}
  \item \textbf{Gate-score} (diagnostic): mask scores at layer $K$; only the selection changes
        $\Rightarrow$ causal test of H1.
  \item \textbf{Gate-input} (deployment): drop background before the LLM; air removal is free
        compression.
  \item Composes with any scorer (FastV, PDrop, DivPrune, Uniform). \textbf{Gate+Uniform} is the
        simplest possible method; matching MedSTP with it would be the headline.
  \item No-op on modalities without background (pathology, dermoscopy) $\Rightarrow$ built-in
        negative control.
  \item \textbf{The ablation identifies the cause:} Gate vs.\ RoPE-free scoring~\cite{feather}
        vs.\ artifact exclusion~\cite{ttreg} vs.\ all combined.
\end{itemize}

% ------------------------------------------------------------------
\section*{7\quad Plan (inference only, $\approx$40 T4-hours on Kaggle)}
\begin{itemize}
  \item \textbf{Models:} Lingshu-7B (M-RoPE) and HuatuoGPT-Vision-7B (1D RoPE), a free
        positional-encoding contrast for H2. Qwen2.5-VL-7B on GQA/POPE as the natural-image
        control.
  \item \textbf{Data:} SLAKE (CT/MRI/X-ray, organ masks), VQA-RAD, OmniMedVQA low-background
        subset (negative control).
  \item \textbf{Baselines:} Full, text-only, Random (3 seeds), Uniform, FastV,
        RoPE-free~\cite{feather}, VisionZip~\cite{visionzip}, DivPrune~\cite{divprune},
        PDrop~\cite{pdrop}; MedSTP quoted from its paper.
  \item \textbf{Metrics:} accuracy on all questions and on the image-dependent subset (full model
        correct, text-only wrong); background and evidence retention; kept-position
        distribution; paired bootstrap CIs and McNemar tests. Kept-token indices are logged, so
        all diagnostics run offline.
  \item \textbf{Timeline:} 2 Nov kill test (SLAKE, 10\%) $\cdot$ 3--9 Nov full grid, diagnostics,
        ablations $\cdot$ 10 Nov abstract registration (deadline 11 Nov).
  \item \textbf{Kill criteria:} (K1) FastV $\ge$ Random on SLAKE CT/MRI in our pipeline
        $\Rightarrow$ stop. (K2) FastV background retention $\le$ Random's $\Rightarrow$ H1
        rejected; test H2/H3. (K3) RoPE-free scoring alone $\ge$ Gate $\Rightarrow$ reframe
        around position.
\end{itemize}

% ------------------------------------------------------------------
\begingroup\small
\begin{thebibliography}{99}
\setlength{\itemsep}{0pt}
\bibitem{fastv} L.~Chen et al. An image is worth 1/2 tokens after layer 2: Plug-and-play
  inference acceleration for large vision-language models. ECCV 2024. arXiv:2403.06764.
\bibitem{medstp} M.~Fan et al. MedSTP: Preserving subspace and transition evidence for medical
  VQA token pruning. arXiv:2610.22751, 2026.
\bibitem{vitos} K.~Chen et al. Token-sparse medical multimodal reasoning via dual-stream
  reinforcement learning. ICML 2026. arXiv:2606.31599.
\bibitem{feather} M.~Endo et al. Feather the throttle: Revisiting visual token pruning for
  vision-language model acceleration. ICCV 2025.
\bibitem{pore} Attention debiasing for token pruning in vision language models.
  arXiv:2508.17807, 2025.
\bibitem{randtok} Y.~Wang et al. When token pruning is worse than random: Understanding visual token
  information in VLLMs. arXiv:2512.07580, 2025.
\bibitem{registers} T.~Darcet, M.~Oquab, J.~Mairal, P.~Bojanowski. Vision transformers need
  registers. ICLR 2024.
\bibitem{ttreg} N.~Jiang et al. Vision transformers don't need trained registers.
  arXiv:2506.08010, 2025.
\bibitem{vtcbench} Are we using the right benchmark: An evaluation framework for visual token
  compression methods. ACL 2026. arXiv:2510.07143.
\bibitem{pdrop} L.~Xing et al. PyramidDrop: Accelerating your large vision-language models via
  pyramid visual redundancy reduction. CVPR 2025.
\bibitem{visionzip} S.~Yang et al. VisionZip: Longer is better but not necessary in vision
  language models. CVPR 2025.
\bibitem{divprune} S.~R.~Alvar et al. DivPrune: Diversity-based visual token pruning for large
  multimodal models. CVPR 2025.
\bibitem{huatuo} J.~Chen et al. HuatuoGPT-Vision: Towards injecting medical visual knowledge
  into multimodal LLMs at scale. arXiv:2406.19280, 2024.
\end{thebibliography}
\endgroup

\end{document}
